% (C) 2026 Ralf Stephan, in collaboration with Claude Code.
% Released under CC0 1.0 Universal (public-domain dedication).
% See https://creativecommons.org/publicdomain/zero/1.0/

\documentclass[11pt]{article}
\usepackage[T1]{fontenc}
\usepackage[utf8]{inputenc}
\usepackage{lmodern}
\usepackage[margin=2.5cm]{geometry}
\usepackage{amsmath,amssymb}
\usepackage{booktabs,tabularx,xltabular}
\usepackage{braket}
\usepackage[hidelinks]{hyperref}

% Lean names and paths: verbatim, breakable at . _ /
\DeclareUrlCommand\lean{\urlstyle{tt}}
\def\UrlBreaks{\do\.\do\_\do\/}
\newcolumntype{L}{>{\raggedright\arraybackslash}X}
\newcolumntype{B}{>{\raggedright\arraybackslash}p{0.42\textwidth}}
\emergencystretch=6em

\title{Lean formalization of an Assumptions of Physics technical brief:\\ log}
\author{}
\date{}

\begin{document}
\maketitle

\noindent
This log records which statements of the technical brief in this directory were checked in Lean,
where the formal version differs from the brief, and any problems found in the brief.
The brief is available at \url{https://assumptionsofphysics.org/resources/briefs}.

\section*{Setup}

\begin{itemize}
\item Lean \texttt{v4.36.0-rc2}; Mathlib at commit \texttt{1b1ccf4b}; physlib (\texttt{master})
  and TauCeti (commit \texttt{e7e7f52}, on the same Mathlib) as further dependencies.
\item Build: \texttt{lake build} in the project root.
\item Every file is checked with \texttt{\#print axioms} on its main results. A result is
  recorded here only if it depends on no axioms beyond \lean{propext},
  \lean{Classical.choice} and \lean{Quot.sound}, and so has no \lean{sorry}.
\item Imports are kept to the Mathlib modules actually needed.
\item Material that belongs in Mathlib is kept in \lean{ForMathlib/}, in Mathlib's directory
  layout (library \lean{ForMathlib}), and imported by the brief files.
\end{itemize}


%======================================================================

\section{Projective measurements are equilibration processes}

\paragraph{Files.}\leavevmode
\begin{itemize}
\item File: \lean{briefs/ProjectiveMeasurement/ProjectiveMeasurement.lean}
\item Library: \lean{ProjectiveMeasurement}
\item Imports:\newline \lean{ForPhyslib.QuantumMechanics.Dephasing},\newline
  \lean{ForPhyslib.QuantumMechanics.RelativeEntropy},\newline
  and the brief \lean{UnitaryEntropy} (for \lean{vonNeumann})
\item New files in \lean{ForPhyslib/QuantumMechanics/}: \lean{Pinching} (spectral projectors
  as the functional calculus of indicators, the pinching, the commutant, trace duality, density
  operators and supports, the Hilbert--Schmidt Pythagoras identity), \lean{Dephasing} (the
  Lindblad generator, the explicit solution, uniqueness, the limit), \lean{RelativeEntropy}
  (relative entropy, Klein's inequality with its equality case, the relative-entropy Pythagoras
  identity). They use from TauCeti:
  \begin{itemize}
  \item the Hilbert--Schmidt norm,
    \lean{TauCeti.Analysis.InnerProductSpace.HilbertSchmidt.Basic};
  \item finiteness of the spectrum, \lean{TauCeti.Analysis.Normed.Operator.FiniteDimensional};
  \item the eigenvector action of the functional calculus,\newline
    \lean{TauCeti.Analysis.InnerProductSpace.FunctionalCalculusComparison};
  \item the supporting-line bound for $x\log x$,
    \lean{TauCeti.Analysis.SpecialFunctions.Log.MulLog}.
  \end{itemize}
\item New \lean{ForMathlib} material:\newline
  \lean{ForMathlib/Analysis/Calculus/MeanValue.lean} (a function continuous from the right at
  $a$ with derivative zero on $(a, \infty)$ is constant on $[a, \infty)$).
\end{itemize}

\paragraph{Setting.} $E$ is a finite-dimensional complex Hilbert space with density operators as
in the brief ``Schr\"odinger equation is equivalent to information conservation''. A process is a
map $\Phi$ sending an observable $X$ and an incoming ensemble $\rho$ to $\Phi(X, \rho)$; the
conditions constrain it for self-adjoint $X$ and density operators $\rho$. For self-adjoint $X$,
\lean{spectralProj}~$X$~$x$ is \lean{cfc} of the indicator of $\{x\}$, the projector $P_x$ onto
the $x$-eigenspace, and \lean{pinching}~$X$~$\rho = \sum_x P_x\rho P_x$, the sum running over the
real spectrum of $X$. \lean{dephasingGenerator}~$X$~$\rho = X\rho X - \tfrac12(X^2\rho +
\rho X^2)$. The Hilbert--Schmidt norm is TauCeti's \lean{hilbertSchmidtNorm}. The relative entropy
\lean{relEntropy}~$\rho$~$\sigma \in [0, \infty]$ is
$\operatorname{tr}(\rho\log\rho) - \operatorname{tr}(\rho\log\sigma)$ if
$\ker\sigma \subseteq \ker\rho$ and $\infty$ otherwise, with $\log = $ \texttt{cfc Real.log}, which
is the logarithm on the support since $\log 0 = 0$. \lean{IsUniqueMinOn}~$f$~$s$~$a$ says that $a
\in s$ and $f(a) < f(b)$ for every other $b \in s$; \lean{commutingDensities}~$X$ is the set of
density operators commuting with $X$. The conditions are
\begin{itemize}
\item \lean{EqMeas}: $\Phi(X, \rho) = P_X(\rho)$;
\item \lean{EqLind}: every $r : \mathbb{R} \to \mathrm{End}(E)$ with $r(0) = \rho$,
  $r'(t) = L(r(t))$ for $t > 0$ and $r$ continuous from the right at $0$ tends to
  $\Phi(X, \rho)$;
\item \lean{EqHS}: $\Phi(X, \rho)$ is the unique minimizer of $\|\rho - \sigma\|_{HS}$ over
  \lean{commutingDensities}~$X$;
\item \lean{EqKL}: $\Phi(X, \rho)$ is the unique minimizer of $D(\rho\|\sigma)$ over the same
  set.
\end{itemize}

\paragraph{Results.}
\begin{center}
\begin{tabularx}{\textwidth}{@{}BL@{}}
\toprule
Brief & Lean \\
\midrule
Eqs.~(3)--(4), blockwise evolution & \lean{spectralProj_mul_dephasingGenerator_mul},
  \lean{dephasingGenerator_block} \\
Eq.~(5), the solution & \lean{dephasing}, \lean{hasDerivAt_dephasing}, \lean{dephasing_zero};
  uniqueness \lean{eq_dephasing} \\
Eqs.~(6)--(10), the limit & \lean{tendsto_dephasing}, \lean{tendsto_pinching_of_hasDerivAt} \\
EQ-LIND $\Leftrightarrow$ EQ-MEAS & \lean{eqMeas_iff_eqLind} \\
``$[\sigma, X] = 0$ iff $P_X(\sigma) = \sigma$'' & \lean{commute_iff_pinching_eq} \\
$P_X$ is positive and trace preserving & \lean{isDensity_pinching} \\
Eq.~(11) & \lean{sq_hilbertSchmidtNorm_sub}, \lean{hilbertSchmidtNorm_sub_pinching_lt} \\
EQ-HS $\Leftrightarrow$ EQ-MEAS & \lean{eqMeas_iff_eqHS} \\
``$\operatorname{tr}(\rho\log\sigma) = \operatorname{tr}(P_X(\rho)\log\sigma)$'' &
  \lean{trace_pinching_mul_of_commute} \\
Eq.~(12) & \lean{relEntropy_eq_add} \\
Eq.~(13) & \lean{relEntropy_pinching}, \lean{relEntropy_pinching_eq} \\
``non-negative and vanishes iff $\sigma = P_X(\rho)$'' (Klein) &
  \lean{re_trace_mul_log_sub_log_nonneg}, \lean{re_trace_mul_log_sub_log_eq_zero_iff},
  \lean{relEntropy_eq_zero_iff} \\
Supports in eq.~(12) & \lean{ker_pinching_le}, \lean{ker_le_ker_pinching} \\
EQ-KL $\Leftrightarrow$ EQ-MEAS & \lean{eqMeas_iff_eqKL}, via \lean{relEntropy_pinching_lt} \\
Theorem~1, complete & \lean{projective_measurement_tfae} \\
\bottomrule
\end{tabularx}
\end{center}

\paragraph{Differences from the brief.}
\begin{itemize}
\item Finite dimension only; the brief has a separable Hilbert space. Then every self-adjoint $X$
  has pure point spectrum, every density operator has finite entropy, so [B]FINENT holds and
  the two halves of Theorem~1 merge into one equivalence of four conditions. The convergence in
  EQ-LIND is in operator norm, which in finite dimension is trace-norm convergence; the
  non-commutative Scheff\'e lemma and dominated convergence are not needed.
\item The Lindblad equation is assumed only for $t > 0$, with continuity from the right at
  $t = 0$. Under these assumptions the solution is unique (\lean{eq_dephasing}), which the brief
  takes for granted when it solves eq.~(4). EQ-LIND is stated for every solution; the explicit
  one shows that there is one.
\item The Hilbert--Schmidt Pythagoras identity and the strict inequality hold for every operator
  $\sigma$ commuting with $X$, not only for density operators; the constraint $\sigma \in D(H)$ in
  EQ-HS is not needed. For EQ-KL it is, through Klein's inequality.
\item Eq.~(13) is stated with the entropy \lean{vonNeumann} of the brief ``Schr\"odinger equation
  is equivalent to information conservation'' and the identity
  $S(\rho) = -\operatorname{tr}(\rho\log\rho)$ (\lean{vonNeumann_eq_neg_re_trace}). As a value in
  $[0, \infty]$ it is $\max(S(P_X(\rho)) - S(\rho), 0)$, which equals $S(P_X(\rho)) - S(\rho)$ by
  Klein's inequality.
\item The relative entropy takes values in $[0, \infty]$, as physlib's \lean{qRelativeEnt}.
  physlib's pinching and \lean{pinching_pythagoras} are stated for matrices and pinch by the
  eigenprojectors of a state; this formalization works with operators on $E$, as the other
  briefs do, and pinches by an observable.
\end{itemize}

\paragraph{Problems found in the brief.} No errors in the statements. Gaps and implicit points:
\begin{itemize}
\item Klein's inequality, ``the last term is non-negative and vanishes if and only if
  $\sigma = P_X(\rho)$'', is used without a reference. It is proved here in eigenbases: with
  $c_{ij} = |\braket{u_i|v_j}|^2$,
  $D(\rho\|\sigma) = \sum_{ij} c_{ij}(r_i\log r_i - r_i\log s_j - r_i + s_j) \ge
  \sum_{ij} c_{ij}(\sqrt{r_i} - \sqrt{s_j})^2$.
\item ``For non-faithful $\sigma$, the identity is understood on its support'' hides two support
  facts that eq.~(12) needs: $\ker P_X(\rho) \subseteq \ker\rho$ for positive $\rho$, and
  $\ker\sigma \subseteq \ker\rho \Rightarrow \ker\sigma \subseteq \ker P_X(\rho)$ for $\sigma$
  commuting with $X$. The first also gives eq.~(12) when $\ker\sigma \not\subseteq \ker\rho$:
  then both sides are infinite.
\item ``Both $\rho(t)$ and $P_X(\rho(0))$ are density operators'': that $\rho(t)$ stays a density
  operator is asserted, not shown. It holds because the matrix
  $\bigl(e^{-t(x-y)^2/2}\bigr)_{x,y}$ is positive semidefinite, so by the Schur product theorem the
  blockwise multiplication preserves positivity. In infinite dimension the Scheff\'e step needs
  it; in finite dimension it is not needed and is not formalized.
\item The sentence ``This base condition states the ensembles have finite entropy'' follows
  QSTO:H, which does not say this; it describes [B]FINENT, stated next.
\item The remark on observables with continuous spectrum is not formalized; in finite dimension
  there is none.
\end{itemize}

\subsection{Hilbert spaces of any dimension}\label{sec:projective-measurement-any-dim}

This subsection covers Theorem~1 for bounded observables with pure point spectrum: EQ-MEAS,
EQ-LIND and EQ-HS are equivalent, and under [B]FINENT EQ-KL is equivalent to EQ-MEAS. It also
covers the claims of the conclusion on repeatability and state update.

\paragraph{Files.}\leavevmode
\begin{itemize}
\item File: \lean{briefs/ProjectiveMeasurement/ProjectiveMeasurement/InfiniteDimensional.lean}
\item Library: \lean{ProjectiveMeasurement}, which now builds every module
  \lean{ProjectiveMeasurement.*}
\item Imports:\newline \lean{ForPhyslib.QuantumMechanics.DephasingFlow},\newline
  \lean{ForPhyslib.QuantumMechanics.VonNeumannRelativeEntropy},\newline
  and the finite-dimensional \lean{ProjectiveMeasurement} (for \lean{IsUniqueMinOn})
\item New file \lean{ForPhyslib/QuantumMechanics/DiscreteObservable.lean}: resolutions of the
  identity and observables with pure point spectrum, the pinching as a strong sum, blocks and
  block-diagonal operators, positivity and trace of the pinching, and the Hilbert--Schmidt
  Pythagoras identity, through TauCeti's block splitting of the Hilbert--Schmidt energy
  (\lean{TauCeti.Analysis.InnerProductSpace.HilbertSchmidt.Pythagoras}).
\item New file \lean{ForPhyslib/QuantumMechanics/DephasingFlow.lean}: the bounded generator and
  its exponential, the blocks of a solution, the unitary group $e^{isX}$, the Gaussian average and
  its trace-norm contraction, equilibration in trace norm.
\item New file \lean{ForPhyslib/QuantumMechanics/VonNeumannRelativeEntropy.lean}: eigenbases,
  the entropy as $\sum_i\eta(\lambda_i)$, the cross term $-\operatorname{tr}(\rho\log\sigma)$,
  Lindblad's formula, Klein's inequality, and the behaviour under pinching.
\item New file \lean{ForMathlib/Analysis/InnerProductSpace/CFCEigenvector.lean}:
  $g(T)v = g(\mu)v$ for $Tv = \mu v$, $T$ self-adjoint and $g$ continuous on the spectrum, in any
  dimension (TauCeti's version is finite-dimensional).
\item Infinite distance when $S(P_X(\rho)) = \infty$, in a module of its own, so that the
  comparator lane is unchanged:\newline
  \lean{briefs/ProjectiveMeasurement/ProjectiveMeasurement/InfiniteDistance.lean}\newline
  Imports: \lean{ProjectiveMeasurement.InfiniteDimensional} and\newline
  \lean{ForPhyslib.QuantumMechanics.VonNeumannEntropyMixture}\newline
  (the latter for a density operator of infinite entropy).
\item Agreement of the two relative entropies, in a module of its own, so that the comparator lane
  is unchanged:\newline
  \lean{briefs/ProjectiveMeasurement/ProjectiveMeasurement/RelativeEntropies.lean}\newline
  Imports: \lean{ProjectiveMeasurement.InfiniteDimensional} and\newline
  \lean{ForPhyslib.QuantumMechanics.RelativeEntropyAgreement}
\item New file \lean{ForPhyslib/QuantumMechanics/RelativeEntropyAgreement.lean}: in finite
  dimension, \lean{relativeEntropy} and \lean{relEntropy} agree on density operators, and every
  entropy is finite.
\item New file:\newline \lean{ForPhyslib/QuantumMechanics/LindbladSumBasisIndependence.lean}\newline
  Lindblad's formula in terms of the eigenspace projections, so that it does not depend on the
  eigenbases.
\item The bridge from the spectral theorem, in a module of its own, so that the comparator lane
  is unchanged:\newline
  \lean{briefs/ProjectiveMeasurement/ProjectiveMeasurement/SpectralTheorem.lean}\newline
  Imports: \lean{ProjectiveMeasurement.InfiniteDimensional} and\newline
  \lean{ForPhyslib.QuantumMechanics.PurePointSpectralMeasure}
\item New file:\newline \lean{ForPhyslib/QuantumMechanics/PurePointSpectrum.lean}\newline
  eigenspace projections \lean{eigenproj}, pure point spectrum, and the existence and uniqueness
  of eigendecompositions.
\item New file:\newline \lean{ForPhyslib/QuantumMechanics/PurePointSpectralMeasure.lean}\newline
  the spectral projections of points are the eigenspace projections, for the spectral measure
  \lean{boundedSelfAdjointSpectralMeasure} of a bounded self-adjoint operator in PhyslibAlpha.
\item Repeatability and state update, in a module of its own, so that the comparator lane is
  unchanged:\newline
  \lean{briefs/ProjectiveMeasurement/ProjectiveMeasurement/Repeatability.lean}\newline
  Imports: \lean{ProjectiveMeasurement.InfiniteDimensional} and\newline
  \lean{ForPhyslib.QuantumMechanics.LuedersRule}
\item New file \lean{ForPhyslib/QuantumMechanics/LuedersRule.lean}: outcome probabilities
  $p_x = \operatorname{tr}(P_x\rho P_x)$, the L\"uders state $\rho_x = P_x\rho P_x/p_x$,
  repeatability, and conditioning of mixtures.
\item EQ-LIND by the route of the brief, eqs.~(6)--(10), in a module of its own, so that the
  comparator lane is unchanged:\newline
  \lean{briefs/ProjectiveMeasurement/ProjectiveMeasurement/HilbertSchmidtRoute.lean}\newline
  Imports: \lean{ProjectiveMeasurement.InfiniteDimensional} and\newline
  \lean{ForPhyslib.QuantumMechanics.DephasingHilbertSchmidt}
\item New file:\newline \lean{ForPhyslib/QuantumMechanics/NoncommutativeScheffe.lean}\newline
  trace-norm bounds
  averaged over the spectral expansion, the gentle measurement bound, and the non-commutative
  Scheff\'e lemma for weak and for Hilbert--Schmidt convergence.
\item New file:\newline \lean{ForPhyslib/QuantumMechanics/DephasingHilbertSchmidt.lean}\newline
  solutions of the dephasing equation stay density operators, their Hilbert--Schmidt distance to
  $P_X(\rho)$, its convergence by dominated convergence (Mathlib's Tannery theorem,
  \lean{Mathlib.Analysis.Normed.Group.Tannery}), and trace-norm convergence through the
  Scheff\'e lemma.
\end{itemize}

\paragraph{Setting.} $E$ is a complex Hilbert space of any dimension; separability is not used.
Density operators, pure states and the entropy \lean{vonNeumannEntropy} $\in [0,\infty]$ are as
in the section ``Hilbert spaces of any dimension'' of
\lean{briefs/SchroedingerEvolution/formalization-log.tex}. An observable is $X \in \mathrm{End}(E)$ with
\lean{IsEigenDecomposition}~$X$~$P$~$x$: the $P_x$ are pairwise orthogonal projections with
$\sum_x P_x = 1$ strongly, the $x$ are distinct reals, and $XP_x = xP_x$; the index set is
arbitrary. The pinching \lean{IsResolutionOfIdentity.pinching}~$\rho = \sum_x P_x\rho P_x$ is a
strong sum, defined for every bounded $\rho$. The conditions are those of finite dimension, with
these changes:
\begin{itemize}
\item \lean{EqLind}: solutions are differentiable in operator norm for $t > 0$ and continuous from
  the right at $0$; they tend to $\Phi(X, \rho)$ in trace norm.
\item \lean{EqHS}: minimization over \lean{commutingDensityOperators}~$X$, with TauCeti's
  \lean{hilbertSchmidtNorm}.
\item \lean{EqKL}: only for $\rho$ with $S(\rho)$ and $S(P_X(\rho))$ finite, and minimization over
  the density operators of finite entropy commuting with $X$. \lean{EqMeasFinEnt} is EQ-MEAS on
  the same $\rho$.
\item The relative entropy \lean{relativeEntropy} is Lindblad's formula
  $\sum_{ij}|\braket{e_i|f_j}|^2(\lambda_i\log\lambda_i - \lambda_i\log\mu_j + \mu_j - \lambda_i)$
  in eigenbases of $\rho$ and $\sigma$ chosen once, each term in $[0,\infty]$ and $+\infty$ for
  $\mu_j = 0 < \lambda_i$.
\end{itemize}

\paragraph{Results.}
\begin{xltabular}{\textwidth}{@{}BL@{}}
\toprule
Brief & Lean \\
\midrule
\endfirsthead
\toprule
Brief & Lean \\
\midrule
\endhead
\bottomrule
\endlastfoot
Eq.~(2), blocks determine $\rho$ & \lean{IsResolutionOfIdentity.ext_block} \\
Eqs.~(3)--(4) & \lean{IsEigenDecomposition.mul_dephasingGenerator_mul} \\
Eq.~(5), for every solution & \lean{IsEigenDecomposition.mul_eq_of_hasDerivAt} \\
A solution exists & \lean{dephasingFlow}, \lean{hasDerivAt_dephasingFlow} \\
$\rho(t)$ is a mixture of unitary conjugations of $\rho(0)$ &
  \lean{gaussianDephasing}, \lean{IsEigenDecomposition.mul_gaussianDephasing_mul},
  \lean{IsEigenDecomposition.eq_gaussianDephasing} \\
Eqs.~(6)--(10), trace-norm convergence &
  \lean{IsEigenDecomposition.traceENorm_gaussianDephasing_sub_le},
  \lean{IsEigenDecomposition.tendsto_traceENorm_of_hasDerivAt} \\
``Both $\rho(t)$ and $P_X(\rho(0))$ are density operators'' &
  \lean{IsEigenDecomposition.gaussianDephasing_nonneg},
  \lean{IsEigenDecomposition.isDensityOperator_of_hasDerivAt} \\
Eqs.~(6)--(9), the brief's route: Hilbert--Schmidt convergence &
  \lean{IsEigenDecomposition.hilbertSchmidtEnergy_sub_pinching_of_hasDerivAt},
  \lean{IsEigenDecomposition.tendsto_hilbertSchmidtNorm_of_hasDerivAt} \\
Non-commutative Scheff\'e lemma, between eqs.~(9) and~(10) &
  \lean{tendsto_traceENorm_of_tendsto_hilbertSchmidtNorm}, for weak convergence
  \lean{tendsto_traceENorm_of_tendsto_inner}; gentle measurement
  \lean{traceENorm_sub_finiteProj_le} \\
Eq.~(10) and EQ-LIND $\Leftrightarrow$ EQ-MEAS, the brief's route &
  \lean{IsEigenDecomposition.tendsto_traceENorm_of_hasDerivAt'}, \lean{eqMeas_iff_eqLind'} \\
EQ-LIND $\Leftrightarrow$ EQ-MEAS & \lean{eqMeas_iff_eqLind} \\
``$[\sigma, X] = 0$ iff block diagonal iff $P_X(\sigma) = \sigma$'' &
  \lean{IsEigenDecomposition.commute_iff}, \lean{IsResolutionOfIdentity.pinching_eq_self_iff} \\
$P_X$ is positive and trace preserving & \lean{IsResolutionOfIdentity.pinching_nonneg},
  \lean{posTrace_pinching}, \lean{isDensityOperator_pinching} \\
Eq.~(11) & \lean{IsResolutionOfIdentity.hilbertSchmidtEnergy_sub},
  \lean{hilbertSchmidtNorm_sub_pinching_lt} \\
EQ-HS $\Leftrightarrow$ EQ-MEAS & \lean{eqMeas_iff_eqHS} \\
$-\operatorname{tr}(\rho\log\sigma)$ & \lean{crossEntropy},
  \lean{IsEigenbasis.crossEntropy_eq} \\
Bridge from the spectral theorem: pure point spectrum gives $X = \sum_x x\,\mu_S(\{x\})$ &
  \lean{IsSelfAdjoint.isEigenDecomposition_eigenproj},
  \lean{isEigenDecomposition_spectralMeasureProj}, \lean{isEigenDecomposition_spectral} \\
$\mu_S(\{x\})$ is the eigenspace projection, and $P_x = \mu_S(\{x\})$ in every
  eigendecomposition & \lean{spectralMeasureProj_singleton},
  \lean{IsEigenDecomposition.eigenproj_eq},
  \lean{IsEigenDecomposition.spectralMeasureProj_singleton_eq} \\
Pure point spectrum iff $\sum_x \mu_S(\{x\}) = 1$ strongly &
  \lean{hasPurePointSpectrum_iff_hasSum_eigenproj},
  \lean{hasPurePointSpectrum_iff_hasSum_spectralMeasureProj},
  \lean{IsEigenDecomposition.hasPurePointSpectrum} \\
$P_X$ depends on $X$ only; EQ-MEAS with $P_x = \mu_S(\{x\})$ &
  \lean{IsEigenDecomposition.pinching_eq_pinching}, \lean{eqMeas_iff_spectral} \\
$D(\rho\|\sigma)$ does not depend on the eigenbases &
  \lean{IsEigenbasis.lindbladSum_eq_lindbladSum},
  \lean{relativeEntropy_eq_lindbladSum_of_isEigenbasis},
  \lean{relativeEntropy_eq_tsum_eigenproj} \\
$D(\rho\|\sigma) = -\operatorname{tr}(\rho\log\sigma) - S(\rho)$ &
  \lean{crossEntropy_eq_lindbladSum_add}, \lean{relativeEntropy_add_vonNeumannEntropy},
  \lean{relativeEntropy_eq_lindbladSum} \\
``$\operatorname{tr}(\rho\log\sigma) = \operatorname{tr}(P_X(\rho)\log\sigma)$'' &
  \lean{IsResolutionOfIdentity.crossEntropy_pinching} \\
Eq.~(12) & \lean{IsResolutionOfIdentity.relativeEntropy_eq_add} \\
Eq.~(13) & \lean{IsResolutionOfIdentity.relativeEntropy_pinching},
  \lean{relativeEntropy_pinching_eq} \\
Klein, ``vanishes iff $\sigma = P_X(\rho)$'' & \lean{eq_of_lindbladSum_eq_zero},
  \lean{relativeEntropy_eq_zero_iff} \\
EQ-KL $\Leftrightarrow$ EQ-MEAS under [B]FINENT & \lean{eqKL_iff_eqMeasFinEnt}, via
  \lean{IsResolutionOfIdentity.relativeEntropy_pinching_lt} \\
Theorem~1 & \lean{projective_measurement_tfae} and \lean{eqKL_iff_eqMeasFinEnt} \\
$D(\rho\|\sigma) \ge S(P_X(\rho)) - S(\rho)$ for $[\sigma, X] = 0$ &
  \lean{IsResolutionOfIdentity.vonNeumannEntropy_pinching_le} \\
``Every ensemble commuting with $X$ is infinitely distant from $\rho$'', for $S(\rho) < \infty$ &
  \lean{relativeEntropy_eq_top_of_vonNeumannEntropy_pinching_eq_top} \\
The same, for $S(\rho) = \infty$: false & \lean{not_forall_relativeEntropy_eq_top},
  \lean{isEigenDecomposition_zero} \\
In finite dimension, the two relative entropies agree & \lean{relativeEntropy_eq_relEntropy},
  \lean{IsDensityOperator.vonNeumannEntropy_ne_top} \\
In finite dimension, EQ-KL is the same minimization in both parts &
  \lean{commutingDensities_eq}, \lean{eqKL_iff_relativeEntropy} \\
Sec.~5, ``all outputs are equilibria'': $L(\sigma) = 0$ iff $[\sigma, X] = 0$, and such $\sigma$ is
  its own output & \lean{dephasingGenerator_eq_zero_iff}, \lean{EqLind.apply_eq_self} \\
Sec.~5, repeatability: $\Phi(X, \Phi(X, \rho)) = \Phi(X, \rho)$, and the outcome probabilities
  are unchanged & \lean{EqLind.apply_apply}, \lean{EqMeas.outcomeProb_apply},
  \lean{IsResolutionOfIdentity.tsum_outcomeProb_eq_one} \\
Sec.~5, after the outcome $x$ the state $\rho_x$ is an equilibrium and a second measurement gives
  $x$ with certainty & \lean{IsResolutionOfIdentity.isDensityOperator_luders},
  \lean{commute_luders}, \lean{IsResolutionOfIdentity.outcomeProb_luders} \\
Sec.~5, ``the output is a classical statistical mixture'': $\Phi(X, \rho) = \sum_x p_x\rho_x$ &
  \lean{IsResolutionOfIdentity.hasSum_pinching_apply_luders}, \lean{EqMeas.hasSum_apply_luders} \\
Sec.~5, ``the update \dots\ is equivalent to what happens in the classical case'' &
  \lean{IsResolutionOfIdentity.luders_of_hasSum}, \lean{EqMeas.luders_apply_eq_component},
  \lean{EqMeas.luders_apply} \\
\end{xltabular}

\paragraph{Differences from the brief.}
\begin{itemize}
\item Observables are bounded with pure point spectrum, given by their eigenprojections and
  distinct eigenvalues. The brief allows any self-adjoint $X$ and remarks that the conditions fail
  on a continuous part of the spectrum; unbounded $X$, where the Lindblad equation holds only on
  a domain, is not formalized. The bridge from the spectral theorem is formalized: a bounded
  self-adjoint $X$ has pure point spectrum (its eigenvectors span a dense subspace) iff its
  spectral measure $\mu_S$ is atomic, $\sum_x\mu_S(\{x\}) = 1$ strongly, and then
  $X = \sum_x x\,\mu_S(\{x\})$. Every eigendecomposition has $P_x = \mu_S(\{x\})$, the
  projection onto $\ker(X - x)$, so the pinching and the conditions depend on $X$ only. The
  identification uses $\|(X - a)v\|^2 = \int|r - a|^2\,d\langle v, \mu_S(r) v\rangle$, through
  PhyslibAlpha's bounded spectral integral.
\item EQ-LIND is proved twice. The first proof (\lean{eqMeas_iff_eqLind}) uses neither the
  Hilbert--Schmidt estimate (6)--(9) nor the non-commutative Scheff\'e lemma. Every solution has
  the blocks of eq.~(5), so it equals the
  Gaussian average $\int e^{isX}\rho\,e^{-isX}\,dN(0,t)(s)$, whose blocks carry the characteristic
  function $e^{-t(x-y)^2/2}$. As an average of unitary conjugations, computed in
  $\mathcal{S}_1$, it is a trace-norm contraction. With $Q = \sum_{x \in F} P_x$ for a finite set
  $F$, $\|\rho(t) - P_X(\rho)\|_1$ is at most $\|\rho - Q\rho Q\|_1 + \|QP_X(\rho)Q - P_X(\rho)\|_1$,
  both small for large $F$ uniformly in $t$, plus finitely many off-diagonal blocks of
  $Q\rho Q$, each decaying like $e^{-t(x-y)^2/2}$.
\item The second proof (\lean{eqMeas_iff_eqLind'}) follows the brief. The Hilbert--Schmidt
  energy of $\rho(t) - P_X(\rho)$ is $\sum_{x \ne y} e^{-t(x-y)^2}\|P_x\rho P_y\|_{HS}^2$, which
  tends to zero by dominated convergence, the bound being $\|\rho\|_{HS}^2 \le 1$. The
  Scheff\'e lemma is proved in the stronger form of Simon, \emph{Trace Ideals}, Thm.~2.20 for
  $p = 1$: density operators converging \emph{weakly} to a density operator converge in trace
  norm. The proof truncates with $Q = \sum_{k \in F}|b_k\rangle\langle b_k|$, the $b_k$
  eigenvectors of the limit $\rho$ carrying weight $> 1 - \eta^2$, and uses the gentle
  measurement bound $\|\sigma - Q\sigma Q\|_1 \le (1 - \operatorname{tr}(Q\sigma))/\eta + \eta$
  (cf.\ Winter 1999, Lemma~9), obtained by averaging the pure-state bound
  $\||c\rangle\langle c| - |Qc\rangle\langle Qc|\|_1 \le 2\|c - Qc\|$ over the spectral expansion
  of $\sigma$. That $\rho(t)$ is a density operator follows from the Gaussian representation:
  the average of positive operators is positive, and its trace is that of its pinching, which
  is $P_X(\rho)$.
\item The relative entropy follows decision~D4: Lindblad's formula, in eigenbases chosen once.
  The cross term $-\operatorname{tr}(\rho\log\sigma)$ is defined without bases, as
  $\sup_n \operatorname{tr}(g_n(\sigma)\rho\,g_n(\sigma))$ for continuous $g_n$ with
  $g_n^2 \uparrow -\log$. In any eigenbases, Lindblad's sum plus $S(\rho)$ equals the cross term,
  so the formula does not depend on the bases once $S(\rho) < \infty$. For every $\rho$ it does
  not depend on them either: with $P_a$ and $Q_b$ the projections onto the eigenspaces of $\rho$
  and $\sigma$, grouping the sum by eigenvalues gives
  $D(\rho\|\sigma) = \sum_{a,b}\operatorname{tr}(Q_bP_aQ_b)\,K(a,b)$ with
  $K(a,b) = a\log a - a\log b + b - a$ ($+\infty$ for $b = 0 < a$), all terms in $[0,\infty]$
  (\lean{relativeEntropy_eq_tsum_eigenproj}). Klein's equality case is read off the sum: a
  vanishing term
  forces $\lambda_i = \mu_j$ or $\braket{e_i|f_j} = 0$, so $\rho$ and $\sigma$ agree on every
  $\braket{e_i|\,\cdot\,f_j}$. In finite dimension it agrees with the finite-dimensional
  \lean{relEntropy} on density operators (\lean{relativeEntropy_eq_relEntropy}): the entropy is
  finite, so Lindblad's formula may be taken in the eigenvector bases, where it is
  $\operatorname{tr}(\rho(\log\rho - \log\sigma))$ if $\ker\sigma \le \ker\rho$ and has a term
  $+\infty$ otherwise. Hence the finite-dimensional EQ-KL minimizes the same function over the
  same set as the part in any dimension (\lean{eqKL_iff_relativeEntropy}).
\item The Pythagoras identity (12) is proved in the form $D(\rho\|\sigma) + S(\rho) =
  -\operatorname{tr}(P_X(\rho)\log\sigma) = D(P_X(\rho)\|\sigma) + S(P_X(\rho))$. The middle step
  holds because each $g_n(\sigma)$ commutes with every $P_x$ and the pinching preserves traces. No
  eigenbasis adapted to the blocks is needed.
\item Under [B]FINENT the hypothesis $S(\rho) < \infty$ is redundant: $S(\rho) \le
  S(P_X(\rho))$ follows from the identities above. \lean{EqKL} keeps it, as the brief does.
\item The Hilbert--Schmidt Pythagoras identity holds in $[0, \infty]$ for every operator, through
  TauCeti's block splitting of the energy, with no summability assumption. The strict inequality
  uses that differences of density operators have finite energy.
\item The two claims of the conclusion are made precise as follows. The outcome $x$ has the
  probability $p_x = \operatorname{tr}(P_x\rho P_x)$, and the state after it is the L\"uders state
  $\rho_x = P_x\rho P_x/p_x$; the brief names neither. Repeatability: the equilibria of the
  dephasing equation are the operators commuting with $X$, under EQ-LIND each commuting density
  operator is its own output, so applying the process to its output returns it; $\rho_x$ is again
  an equilibrium, and gives $x$ with probability one. Classical update: a mixture
  $\sigma = \sum_x w_x\sigma_x$ of density operators $\sigma_x = P_x\sigma_xP_x$ gives the outcome
  $x$ with probability $w_x$, and its L\"uders state for $x$ is $\sigma_x$. This is conditioning
  of a classical mixture on the component that occurred. The output $\sum_x p_x\rho_x$ is such a
  mixture, and its update on $x$ equals the L\"uders state of the input $\rho$.
\end{itemize}

\paragraph{Problems found in the brief.} No errors in the statements. Gaps beyond those of the
finite-dimensional part:
\begin{itemize}
\item ``Both $\rho(t)$ and $P_X(\rho(0))$ are density operators'' is asserted. In infinite
  dimension it is needed for the Scheff\'e step. The Gaussian representation above proves it
  (\lean{IsEigenDecomposition.isDensityOperator_of_hasDerivAt}).
\item The Scheff\'e step itself (``Hilbert--Schmidt convergence to a density operator implies
  trace-norm convergence'') is cited without reference or proof. It is a known theorem (Simon,
  \emph{Trace Ideals}, Thm.~2.20; Dell'Antonio 1967), proved here in its weak-convergence form
  (\lean{tendsto_traceENorm_of_tendsto_inner}).
\item The conclusion's claims on repeatability and state update use the outcome probabilities and
  the post-measurement state without defining them; see the last item of the differences above.
\item The brief solves the block equations (4) but does not say that the solution is unique; the
  argument behind \lean{IsEigenDecomposition.mul_eq_of_hasDerivAt} gives uniqueness for every
  solution continuous from the right at $0$.
\item Eq.~(13)'s ``finite if and only if $S(P_X(\rho))$ is finite'' uses $S(\rho) < \infty$,
  which the admissible domain provides; outside it the equivalence fails, since a block-diagonal
  $\rho$ with $S(\rho) = \infty$ has $D(\rho\|P_X(\rho)) = D(\rho\|\rho) = 0$. The remark
  ``Relative entropy in infinite dimensions'' says that in the case $S(P_X(\rho)) = +\infty$
  ``every ensemble commuting with $X$ is infinitely distant from $\rho$''. Read literally, the same
  $\rho$ is a counterexample, with $\sigma = \rho$: in every infinite-dimensional space, for
  $X = 0$ with the single eigenprojection $1$ and a density operator of infinite entropy
  (\lean{not_forall_relativeEntropy_eq_top}). The claim holds when $S(\rho) < \infty$, by
  $D(\rho\|\sigma) + S(\rho) = -\operatorname{tr}(P_X(\rho)\log\sigma) \ge S(P_X(\rho))$
  (\lean{relativeEntropy_eq_top_of_vonNeumannEntropy_pinching_eq_top}).
\end{itemize}

\end{document}
